\chapter{Coherent Hamiltonian actions}
\label{mht:ch:coherent-actions}

Hamiltonian vector fields act on the commutative plane by a strict
Lie action. Replacing functions by noncommutative words changes the
problem. Symmetrized cyclic derivatives still preserve the
commutativity relation, but their brackets acquire a defect.
The variable resolving that relation supplies the first homotopy.
A contraction of the same resolution supplies every higher one.

The construction also works for the Weyl relation. We keep its
coefficient parameter separate from quantization of matrix entries.
After constructing the actions, we ask which homotopies preserve
the volume and which are Hamiltonian functions on the matrix
gauge complex. These are additional structures on an action.

\section{The first bracket defect}

Work over $\mathbb C$, with differentials of cohomological degree
one. Write $q,p$ for the coordinates of the symplectic plane, and put
\[
 A=\mathbb C[q,p],\qquad
 \mathfrak l=A/\mathbb C1,\qquad
 \{f,g\}=f_qg_p-f_pg_q.
\]
Thus $X_f=\{f,-\}$ is independent of the constant term of $f$.
For a monomial $q^ap^b$, let $\sigma(q^ap^b)$ be the average of
all distinct words with $a$ factors $q$ and $b$ factors $p$.
For a word $w=a_1\cdots a_m$, define
\[
 \partial_q^{\mathrm{cyc}}w
 =\sum_{a_i=q}a_{i+1}\cdots a_ma_1\cdots a_{i-1},
\]
and define $\partial_p^{\mathrm{cyc}}$ in the same way.

Introduce the free graded associative algebra
\begin{equation}\label{mht:eq:relation-resolution}
 R=\mathbb C\langle q,p,t\rangle,
 \qquad |q|=|p|=0,\quad |t|=-1,
 \qquad dq=dp=0,\quad dt=qp-pq.
\end{equation}
Here $t$ is a relation generator. For homogeneous words, brackets
denote graded commutators:
$[a,b]=ab-(-1)^{|a||b|}ba$. The projection $\pi:R\to A$ kills
$t$ and makes $q,p$ commute.

The cyclic lift is the degree-zero derivation
\begin{equation}\label{mht:eq:cyclic-first-action}
 F_1(f)q=-\partial_p^{\mathrm{cyc}}\sigma f,
 \qquad F_1(f)p=\partial_q^{\mathrm{cyc}}\sigma f,
 \qquad F_1(f)t=0.
\end{equation}
It kills constants. Telescoping the rotations of a word gives
\[
 [q,\partial_q^{\mathrm{cyc}}w]
 +[p,\partial_p^{\mathrm{cyc}}w]=0.
\]
Consequently $[d,F_1(f)]=0$ and $\pi F_1(f)=X_f\pi$.

Put $\mu=qp-pq$. Cyclic differentiation of the cubic pair gives
\begin{equation}\label{mht:eq:cubic-action-defect}
 \bigl([F_1(q^3),F_1(p^3)]-F_1(9q^2p^2)\bigr)(q,p,t)
 =\bigl(3[\mu,q],3[\mu,p],0\bigr).
\end{equation}
This defect is nonzero. Substitution of
$X=\operatorname{diag}(1,0)$ and $Y=E_{12}$ for $q,p$ gives
$-3E_{12}$ in its first component.

The defect is exact as a derivation. The degree-minus-one derivation
\begin{equation}\label{mht:eq:cubic-action-homotopy}
 \eta(q,p,t)=\bigl(3[t,q],3[t,p],3t^2\bigr)
\end{equation}
satisfies $[d,\eta]$ equal to \eqref{mht:eq:cubic-action-defect}.
On $q,p$ this follows from $dt=\mu$. On $t$, the terms cancel
because $d(t^2)=\mu t-t\mu$ and
$\eta(\mu)=3[t,\mu]$. The value on $t$ is necessary.

Changing the polynomial cyclic lift cannot make it a Lie section.
Let $\mathcal N$ be the vector space of cyclic words in $q,p$,
modulo constants. Equivalently, quotient the free algebra by the
span of commutators and then by its constant class. Its necklace
bracket is
\[
 \{[U],[V]\}_{\mathcal N}
 =\bigl[\partial_q^{\mathrm{cyc}}U\,
          \partial_p^{\mathrm{cyc}}V
       -\partial_p^{\mathrm{cyc}}U\,
          \partial_q^{\mathrm{cyc}}V\bigr].
\]

Expanding the derivatives identifies its Hamiltonian commutator.
The cyclic Euler identity
$[q\partial_q^{\mathrm{cyc}}U+p\partial_p^{\mathrm{cyc}}U]
=m[U]$ in degree $m$ makes this action faithful modulo constants.
Jacobi therefore follows from Jacobi for derivations.
Abelianization is a Lie map $\mathcal N\to\mathfrak l$.

\begin{proposition}\label{mht:prop:necklace-no-section}
Abelianization $\mathcal N\to\mathfrak l$ has no linear Lie
section, even without a degree-preservation requirement.
\end{proposition}

\begin{proof}
Abelianization is injective in degrees at most three. A lift of a
cubic therefore differs from its unique cyclic cubic by kernel
terms of degree at least four. Their brackets with cubics have
degree at least five.

The cyclic classes $C=[q^2p^2]$ and $D=[qpqp]$ are distinct, and
\[
 \{[q^3],[p^3]\}_{\mathcal N}=9C,
 \qquad
 \{[q^2p],[qp^2]\}_{\mathcal N}=C+2D.
\]
The commutative brackets are $9q^2p^2$ and $3q^2p^2$.
Thus the degree-four part of a Lie lift of $q^2p^2$ would be both
$C$ and $(C+2D)/3$. This would identify the two distinct cyclic
words.
\end{proof}

The proposition concerns sections through cyclic words. To
construct an action on $R$, we instead solve the bracket defects
in the complex of derivations.

\section{A word contraction for the plane and its deformation}

Introduce a degree-zero parameter $\varepsilon$ and set
$k_\varepsilon=\mathbb C[\varepsilon]$. The family
\begin{equation}\label{mht:eq:weyl-relation-resolution}
 \mathcal R_\varepsilon
   =k_\varepsilon\langle q,p,t\rangle,
 \qquad dt=qp-pq-\varepsilon,
 \qquad
 W_\varepsilon
   =k_\varepsilon\langle q,p\rangle/(qp-pq-\varepsilon)
\end{equation}
has the same generator degrees as $R$. Write $\pi_\varepsilon$
for the projection killing $t$. The rule $qp\mapsto pq+\varepsilon$
has no overlapping leading words. It gives the normal basis
$p^bq^a$ of $W_\varepsilon$ and a $k_\varepsilon$-linear section
$i$ of $\pi_\varepsilon$.

This section is a chain map, not an algebra map:
\[
 i(q)i(p)-i(qp)=qp-(pq+\varepsilon)=dt.
\]
Here $qp$ inside $i$ denotes the product in $W_\varepsilon$.
We now construct a homotopy for this linear section.

First replace $dt$ by $d_0t=qp$. Expand each $t$ in a word into
$qp$, marking the occurrence expanded. Occurrences of $qp$ do
not overlap. For each expanded word, the subsets of marked
occurrences form a Boolean complex. The differential $d_0$
removes a mark, with sign $(-1)^a$ when $a$ marks precede it.

Define $h_0$ by adding a mark at the leftmost occurrence if it is
unmarked. Give zero if that occurrence is already marked, or if
there is no occurrence. Let $p_0$ retain only the unmarked normal
words. Removing and adding the distinguished mark gives
\[
 d_0h_0+h_0d_0=1-ip_0,
 \qquad h_0^2=h_0i=p_0h_0=0.
\]
Every removal at another occurrence cancels with the corresponding
addition, with the opposite tensor sign.

Write $d=d_0+\delta$, where $\delta t=-pq-\varepsilon$. The
exchange term lowers the number of pairs with a $q$ before a $p$.
The scalar term lowers expanded-word length. Ordered first by
length and then by this inversion number, every term of
$\delta h_0$ decreases. Hence the sums
\begin{equation}\label{mht:eq:finite-word-contraction}
 h=h_0\sum_{j\geq0}(-\delta h_0)^j,
 \qquad p'=p_0\sum_{j\geq0}(-\delta h_0)^j
\end{equation}
terminate on every polynomial word.

\begin{proposition}\label{mht:prop:word-contraction}
The map $p'$ equals $\pi_\varepsilon$, and
\[
 \pi_\varepsilon i=1,\qquad
 dh+hd=1-i\pi_\varepsilon,
 \qquad h^2=hi=\pi_\varepsilon h=0.
\]
Thus $\mathcal R_\varepsilon$ resolves $W_\varepsilon$ over
$k_\varepsilon$. The same identities hold continuously after
$\varepsilon$-adic completion.
\end{proposition}

\begin{proof}
Insert the finite inverse $(1+\delta h_0)^{-1}$ into the Boolean
contraction. Expanding its product with $1+\delta h_0$ cancels
the perturbation terms. Since $\delta i=0$, the inclusion and
the zero differential on the normal words stay unchanged. This
gives the homotopy identity and its side conditions.

On degree-zero words, $p'$ performs exactly the terminating rule
$qp\mapsto pq+\varepsilon$. Its result is their image in the
normal basis of $W_\varepsilon$. Negative-degree words have zero
projection. Thus $p'=\pi_\varepsilon$.

Every operation has nonnegative $\varepsilon$-order. Modulo
$\varepsilon^m$, only finitely many coefficients of an input
contribute, each with finite polynomial support. The formulas
therefore define compatible continuous maps on all these
quotients. Their identities pass to the separated inverse limit.
\end{proof}

The completed carrier is
$\widehat{\mathcal R}_\varepsilon
=\varprojlim_m\mathcal R_\varepsilon/\varepsilon^m$.
It permits polynomial word support to grow with parameter order.
At $\varepsilon=0$, the proposition gives a contraction of
$R\to A$. Both the polynomial family and its classical fiber
retain this explicit homotopy.

\section{Contracting the derivation kernel}

Give $\operatorname{Der}_{k_\varepsilon}(\mathcal R_\varepsilon)$
the graded commutator and differential $\partial D=[d,D]$.
Let $J_\varepsilon=\ker\pi_\varepsilon$, and define
\[
 K_\varepsilon
 =\{D:\pi_\varepsilon D(q)
          =\pi_\varepsilon D(p)
          =\pi_\varepsilon D(t)=0\}.
\]
These are derivations with all values in $J_\varepsilon$.
For a degree-$j$ derivation, write
$U=Dq$, $V=Dp$, $W=Dt$, and put
$L(U,V)=[U,p]+[q,V]$. Then
\begin{equation}\label{mht:eq:derivation-kernel-differential}
 (\partial D)(q,p,t)
 =\bigl(dU,dV,dW-(-1)^jL(U,V)\bigr).
\end{equation}

\begin{lemma}\label{mht:lem:derivation-contraction}
The operator of degree minus one defined by
\begin{equation}\label{mht:eq:derivation-kernel-homotopy}
 (HD)(q,p,t)
 =\bigl(hU,hV,h(W-(-1)^jL(hU,hV))\bigr)
\end{equation}
preserves $K_\varepsilon$ and satisfies
$\partial H+H\partial=1$ there. It has finite polynomial values
and extends continuously in $\varepsilon$.
\end{lemma}

\begin{proof}
The ideal $J_\varepsilon$ contains every argument of $h$ in the
formula, and $\pi_\varepsilon h=0$. On $q,p$, the identity is
$dh+hd=1$ on $J_\varepsilon$. Since $q,p$ are closed,
$dL(U,V)=L(dU,dV)$. Substitution in the third component leaves
$(dh+hd)W$ and the additional term
\[
 (-1)^j hL\bigl((dh+hd-1)U,(dh+hd-1)V\bigr).
\]

The first expression equals $W$ and the second vanishes. This
proves the identity on arbitrary kernel derivations, including
those which are not closed. Finiteness and continuity follow
from Proposition~\ref{mht:prop:word-contraction}.
\end{proof}

\section{Solving every action equation}

The Moyal product on $k_\varepsilon[q,p]$ is
\[
 f\star g
 =m\exp\!\left(\frac{\varepsilon}{2}
   (\partial_q\otimes\partial_p-\partial_p\otimes\partial_q)
   \right)(f\otimes g).
\]
Its divided commutator
$B_\varepsilon(f,g)=(f\star g-g\star f)/\varepsilon$
is polynomial. This quotient belongs to $k_\varepsilon[q,p]$.
Quotienting
constant polynomials defines the Lie algebra
$\mathfrak l_\varepsilon=k_\varepsilon[q,p]/k_\varepsilon1$.

Weyl symmetrization identifies this product with $W_\varepsilon$.
One can check the identification on auxiliary exponential series:
the product of $\exp(sq+rp)$ and $\exp(s'q+r'p)$ acquires the
factor $\exp(\varepsilon(sr'-rs')/2)$. The relation
$[q,p]=\varepsilon$ gives this identity by differentiation in
the auxiliary variables. Comparing their coefficients gives
the polynomial identity.

Use \eqref{mht:eq:cyclic-first-action} for $F_1$ over
$k_\varepsilon$. The cyclic identity still proves $[d,F_1]=0$.
Cyclic differentiation of a symmetrized polynomial is the
symmetrization of its ordinary derivative. Thus its projection
is the strict derivation $B_\varepsilon(f,-)$ of $W_\varepsilon$.

We specify the signs of the higher equations by a coalgebra.
For a differential graded Lie algebra $\mathfrak g$, let
$s:\mathfrak g\to\mathfrak g[1]$ have degree $-1$.
The codifferential of the cofree conilpotent symmetric coalgebra
$S^c(\mathfrak g[1])$ has corestrictions
\[
 b_1(sa)=-s\partial a,
 \qquad b_2(sa\,sb)=(-1)^{|a|}s[a,b].
\]
An $L_\infty$ morphism is a coaugmented coalgebra map commuting
with these codifferentials. Its component on $n$ shifted
degree-zero inputs is written $sF_n$. The unshifted alternating
map $F_n$ has degree $1-n$.

Suppose the components below arity $n$ have been chosen. Extend
them by the set-partition formula to a coalgebra map
$\mathcal F$, temporarily setting $F_n=0$. Let $\Omega_n$
be the desuspended arity-$n$ corestriction of
$b\mathcal F-\mathcal F b$. Set
\begin{equation}\label{mht:eq:all-arity-recursion}
 F_n=H\Omega_n\qquad(n\geq2).
\end{equation}

The first two residuals are
\begin{align*}
 \Omega_2(f,g)&=[F_1(f),F_1(g)]-F_1(B_\varepsilon(f,g)),\\
 \Omega_3(f,g,k)&=\sum_{\mathrm{cyc}}
   \bigl([F_1(f),F_2(g,k)]-F_2(B_\varepsilon(f,g),k)\bigr).
\end{align*}
At arity four the residual includes brackets of two binary
components. The coalgebra definition fixes these terms and their
signs along with all subsequent ones.

\begin{theorem}\label{mht:thm:coherent-hamiltonian-action}
The recursion defines an $L_\infty$ morphism
\[
 F:\mathfrak l_\varepsilon
   \longrightarrow
   \operatorname{Der}_{k_\varepsilon}(\mathcal R_\varepsilon)
\]
with first component \eqref{mht:eq:cyclic-first-action}.
Every fixed-arity value on polynomial inputs is polynomial.
The maps extend continuously to the parameter completions.
Specialization at $\varepsilon=0$ gives an action of
$\mathfrak l$ on $R$.
\end{theorem}

\begin{proof}
The arity-one identity is $[d,F_1]=0$. Suppose the equations
hold below arity $n$. The square-zero identities for the two
codifferentials imply
\[
 b(b\mathcal F-\mathcal F b)
 +(b\mathcal F-\mathcal F b)b=0.
\]
Its first nonzero corestriction says $\partial\Omega_n=0$.

The binary residual projects to zero because the action on
$W_\varepsilon$ is strict. Moreover, $F_1$ preserves
$J_\varepsilon$ and normalizes $K_\varepsilon$.
Higher components already constructed lie in $K_\varepsilon$.
Every later residual therefore lies in $K_\varepsilon^{2-n}$.

Lemma~\ref{mht:lem:derivation-contraction} gives
$\partial F_n=\Omega_n$. The term $-s\partial F_n$ from $b_1$
cancels the residual. This proves the equations inductively.
There are finitely many partitions at each arity, and the word
contraction terminates. All operations preserve nonnegative
parameter order, so the same formulas define the completed action
coefficient by coefficient.
\end{proof}

Several values make the higher structure visible. At
$\varepsilon=0$, the chosen $F_2(q^3,p^3)$ is precisely
$\eta$ from \eqref{mht:eq:cubic-action-homotopy}. For the triple
$(q^4,p^4,q^2p)$, the recursion gives
\begin{equation}\label{mht:eq:quartic-ternary-action}
 \begin{aligned}
 5F_3q&=-128tq^2t+16tqtq+112t^2q^2,\\
 5F_3p&=-128ptqt-112pt^2q,\\
 5F_3t&=-64t^2qt-56t^3q.
 \end{aligned}
\end{equation}
Substitution of $dt=qp-pq$ gives $\partial F_3=\Omega_3$
on all three generators. The triple $(q^3,p^3,q^2p)$ instead
has zero residual and zero third component.

For $\mathbf f=(q^3,p^3,qp^3)$ and
$\mathbf g=(q^3,p^3,qp^3,q^2p^2)$, another pair of values is
\begin{equation}\label{mht:eq:third-fourth-action}
 \begin{aligned}
 F_3(\mathbf f)q&=-\tfrac92(tpt+pt^2),\\
 F_4(\mathbf g)q&=\tfrac{33}{5}t^3p+\tfrac{39}{5}t^2pt
                  +\tfrac{69}{5}tpt^2+\tfrac{99}{5}pt^3.
 \end{aligned}
\end{equation}
Both components vanish on $p,t$. Applying the differential to
these words gives their respective coalgebra residuals, including
the binary--binary terms at arity four. Thus the specified action
does not truncate after its second or third component.

The deformation has its own first correction:
\[
 B_\varepsilon(q^3,p^4)=12q^2p^3+6\varepsilon^2p.
\]
For this pair the recursion gives
\begin{equation}\label{mht:eq:quantum-binary-action}
 \begin{aligned}
 F_2q&=6(-pqt+ptq-qtp+tpq-\varepsilon t),\\
 F_2p&=6(-p^2t+tp^2),\\
 F_2t&=6(pt^2+tpt).
 \end{aligned}
\end{equation}
For $(q^3,p^4,q^2p)$ it gives
\[
 F_3q=-12tqt+12t^2q,\qquad
 F_3p=-12pt^2,\qquad F_3t=-6t^3.
\]
These identities follow by differentiating the displayed words
with $dt=qp-pq-\varepsilon$. In particular, the term
$-6\varepsilon t$ contributes $6\varepsilon^2$ to the first
binary equation. Omitting it loses the derivative of the Moyal
correction.

\section{Choices of lifts and pointed completion}

Return to $\varepsilon=0$. Let $\operatorname{Der}_\pi(R)$
be the derivations preserving $J=\ker\pi$. They induce
derivations of $A$, giving an exact sequence of complexes
\[
 0\longrightarrow K\longrightarrow\operatorname{Der}_\pi(R)
   \xrightarrow{\pi_*}\operatorname{Der}(A)\longrightarrow0.
\]
Here $K=K_0$ is understood in the classical construction.

Surjectivity is degreewise: lift the values of a derivation on
$q,p$ and set its value on $t$ to zero. This preserves $J$,
although the chosen lift need not commute with $d$.
The contracted kernel makes $\pi_*$ a quasi-isomorphism.

To specify a space of choices, use the complete convolution
Lie algebra
\[
 \mathcal G(L)=\prod_{n\geq1}
   \operatorname{Hom}\bigl(S^n(\mathfrak l[1]),L\bigr),
 \qquad
 \mathcal F^r\mathcal G(L)=\prod_{n\geq r}
   \operatorname{Hom}\bigl(S^n(\mathfrak l[1]),L\bigr).
\]
Its differential uses that of $L$ and the source Chevalley
codifferential. Its bracket uses the reduced coproduct with
the tensor Koszul sign. Degree-one Maurer--Cartan elements are
the $L_\infty$ morphisms in the preceding convention.

Define its simplicial Maurer--Cartan set by tensoring each
arity quotient with polynomial differential forms on a simplex,
then taking the inverse limit in arity. The derived space of
lifts is the homotopy fiber over $f\mapsto X_f$ of the map
induced by $\pi_*$. This fixes the commutative action.

\begin{proposition}\label{mht:prop:contractible-action-lifts}
The derived space of lifts of $X$ to
$\operatorname{Der}_\pi(R)$ is nonempty and contractible.
\end{proposition}

\begin{proof}
The recursion supplies a lift. The map of convolution complexes
is degreewise onto, and its arity-$n$ associated-graded kernel is
\[
 \operatorname{Hom}\bigl(S^n(\mathfrak l[1]),K\bigr).
\]
Postcomposition with $H$ contracts it.

Consider a compatible boundary or horn of polynomial simplex
forms and a lift already chosen in lower arities. Polynomial
forms extend from that boundary or horn. In barycentric
coordinates the ideal of face $i$ is $(x_i,dx_i)$.
Successive face extensions preserve the restrictions already
fixed. Choose a tentative next component with the prescribed
boundary and projection.

Its Maurer--Cartan defect is a closed kernel element tensored
with relative simplex forms. Closedness is the first nonzero
arity of the Bianchi identity. The tensor extension of $H$
contracts this relative kernel. Subtracting its primitive
corrects the component without changing its boundary or
projection. Arity completeness gives all the compatible
corrections. The resulting simplicial map has lifts for all
boundaries, so its nonempty fiber is contractible.
\end{proof}

This fiber imposes neither a symplectic form nor Hamiltonian
functions on the homotopies. Its completion in the formal
coordinates has an additional continuity requirement.
Give $q,p$ weight one and $t$ weight two, and write
$R_w$ for the homogeneous weight space. Define
$\widehat R=\prod_{w\geq0}R_w$, with discrete homogeneous
pieces and their product topology.

\begin{proposition}\label{mht:prop:pointed-action-completion}
For homogeneous classical inputs of polynomial degrees
$d_1,\ldots,d_n$, the derivation $F_n$ shifts weight by
$\sum_i(d_i-2)$. At every fixed arity the action therefore
extends continuously to
\[
 \widehat{\mathfrak l}_{\mathrm{pt}}
 =\prod_{d\geq2}\operatorname{Sym}^d(\mathbb Cq\oplus\mathbb Cp)
\]
acting on $\widehat R$.
\end{proposition}

\begin{proof}
The differential and word contraction preserve weight.
Cyclic differentiation gives the asserted shift for $F_1$.
Brackets add shifts, and $H$ preserves them. Induction proves
the formula for every $F_n$.

For pointed inputs, each $d_i-2$ is nonnegative. At fixed
arity and fixed output weight, only finitely many input
degrees and source weights contribute. Each homogeneous
piece is finite dimensional. The formulas and identities
therefore extend coefficient by coefficient and are continuous.
\end{proof}

This fixed-arity extension gives no analytic summation over
all arities. Its pointed classical source must
also be distinguished from the quantum one. The term
$6\varepsilon^2p$ above leaves the subspace of polynomials
whose ordinary degree is at least two.

\section{Matrix actions and Hamiltonian functions}

For $N\geq1$, the commutative differential graded algebra of
derived representations of \eqref{mht:eq:weyl-relation-resolution} is
\begin{equation}\label{mht:eq:derived-matrix-action-carrier}
 \mathcal R_{N,\varepsilon}
 =\mathbb C[\varepsilon,X_{ij},Y_{ij}]
       \otimes\Lambda(\Psi_{ij}),
 \qquad |\Psi_{ij}|=-1,
 \qquad Q\Psi=[X,Y]-\varepsilon I_N,
\end{equation}
with $QX=QY=0$. Matrix multiplication uses graded commutative
coefficients. Substitution of $X,Y,\Psi$ in a derivation's
generator values defines a derivation of their entries.
The graded Leibniz rule makes substitution preserve brackets
and differentials.

It follows that $F$ defines a $GL_N$-equivariant $L_\infty$
action on $\mathcal R_{N,\varepsilon}$ and its invariant
algebra. At $\varepsilon=0$, giving $X,Y$ weight one and
$\Psi$ weight two also gives the continuous pointed action
on the weight completion.

The matrix complex is not contracted by the universal word
homotopy. At $\varepsilon=0$, the cycle
$\operatorname{Tr}\Psi$ represents a nonzero class in degree
$-1$. Evaluation at $X=Y=0$ maps it to the same nonzero element
of $\Lambda(\Psi_{ij})$ with zero differential. It cannot be
a boundary. Keeping the relation complex is therefore essential.

To include gauge ghosts classically, introduce an odd matrix
$c$ of degree one and put
\[
 \mathcal B_N
 =\mathbb C[X_{ij},Y_{ij}]
       \otimes\Lambda(c_{ij},\Psi_{ij}).
\]

Its even Poisson bracket is specified by
\[
 \{X_{ij},Y_{kl}\}=\delta_{il}\delta_{jk},\qquad
 \{c_{ij},\Psi_{kl}\}=\{\Psi_{kl},c_{ij}\}
       =\delta_{il}\delta_{jk}.
\]
All other generator brackets vanish. Set
$\Theta=\operatorname{Tr}(c[X,Y]-\Psi c^2)$ and
$Q=\{\Theta,-\}$. Then
\begin{equation}\label{mht:eq:action-bfv-differential}
 QX=-[c,X],\quad QY=-[c,Y],\quad Qc=-c^2,
 \qquad Q\Psi=[X,Y]-[c,\Psi].
\end{equation}

Here $[c,\Psi]=c\Psi+\Psi c$. Expanding on the four
generators, with the graded Leibniz rule, gives $Q^2=0$.
These are the Chevalley--Eilenberg gauge differential and
the moment Koszul differential on one algebra.

Every matrix-word derivation extends by $Dc=0$. For a
degree-$j$ derivation and a homogeneous matrix word $M$,
\[
 D[c,M]=(-1)^j[c,DM].
\]
Thus extension commutes with the gauge part of $Q$ in the
graded sense. It preserves the differential and bracket,
and extends the classical $L_\infty$ action to $\mathcal B_N$.

The first action is Hamiltonian. For $f\in A$, put
$J_N(f)=\operatorname{Tr}\sigma f(X,Y)$ and
$H_{J_N(f)}=\{J_N(f),-\}$. Cyclic differentiation gives
$H_{J_N(f)}=F_1(f)$, and conjugation invariance gives
$QJ_N(f)=0$. For the cubic pair, a Hamiltonian primitive is
\begin{equation}\label{mht:eq:hamiltonian-cubic-primitive}
 K=-\tfrac32\operatorname{Tr}(\Psi\mu),\qquad
 QK=-\tfrac32\operatorname{Tr}\mu^2
    =\{J_N(q^3),J_N(p^3)\}-J_N(9q^2p^2),
 \qquad \mu=[X,Y].
\end{equation}
The gauge terms cancel under the graded trace. Among the six
words in $\sigma(q^2p^2)$, four have trace
$\operatorname{Tr}(X^2Y^2)$ and two have trace
$\operatorname{Tr}(XYXY)$, which proves the last equality.

This Hamiltonian homotopy has components
\[
 H_KX=\tfrac32[\Psi,X],\quad
 H_KY=\tfrac32[\Psi,Y],\quad
 H_Kc=-\tfrac32\mu,\quad H_K\Psi=0.
\]
The extension of \eqref{mht:eq:cubic-action-homotopy} instead has
\[
 DX=3[\Psi,X],\quad DY=3[\Psi,Y],\quad
 Dc=0,\quad D\Psi=3\Psi^2.
\]
Both have the same differential, the cubic action defect.
For $N\geq2$, $D$ is not Hamiltonian. Indeed $Dc=0$ would
make a Hamiltonian independent of $\Psi$. Such a polynomial
has no component of degree $-1$, while $D$ is a nonzero
derivation of that degree.

\begin{proposition}\label{mht:prop:central-action-obstruction}
There is no $L_\infty$ morphism from $\mathfrak l$ to
$(\mathcal B_N,Q,\{-,-\})$ whose first component is $J_N$
on the representatives with zero constant term. The same
obstruction survives exact changes of these closed charges.
\end{proposition}

\begin{proof}
The classes of $q,p$ have zero bracket in $\mathfrak l$, but
$\{\operatorname{Tr}X,\operatorname{Tr}Y\}=N$.
The binary morphism equation would make $N$ a $Q$-boundary.
Evaluation at $X=Y=c=\Psi=0$ is a chain map to $\mathbb C$,
since all values in \eqref{mht:eq:action-bfv-differential}
have positive polynomial degree. It sends that proposed
boundary to the nonzero scalar $N$. Exact changes do not
change the induced bracket on cohomology.
\end{proof}

Retaining the central extension $A\to\mathfrak l$ with
$J_N(1)=N$ removes this particular source obstruction.
Choosing representatives vanishing at the origin gives its
scalar cocycle $c(f,g)=\{f,g\}(0)$. This cocycle vanishes on
$(q,p)^2$. Neither change by itself constructs the remaining
Hamiltonian homotopies.

\section{The volume in the normalized mixed complex}

Return first to the classical resolution $R$. Its normalized
Hochschild chains are
$C_n(R)=R\otimes\overline R^{\otimes n}$, where
$\overline R=R/\mathbb C1$. Write a chain as
$a_0[a_1|\cdots|a_n]$. Its internal degree is
$m=\sum_i|a_i|$, and its total cohomological degree is $m-n$.

Use the raw tensor convention
\begin{equation}\label{mht:eq:raw-mixed-convention}
 d_{\mathrm H}=(-1)^n d_\otimes+b.
\end{equation}
The tensor differential $d_\otimes$ has its usual Koszul
signs. The other operators are
\begin{align*}
 b(a_0[a_1|\cdots|a_n])
 ={}&\sum_{j=0}^{n-1}(-1)^j
       a_0[\cdots|a_ja_{j+1}|\cdots|a_n]\\
 &+(-1)^{n+|a_n|\sum_{i<n}|a_i|}
       a_na_0[a_1|\cdots|a_{n-1}],\\
 B(a_0[a_1|\cdots|a_n])
 ={}&\sum_{j=0}^n
   (-1)^{nj+(\sum_{i<j}|a_i|)(\sum_{i\geq j}|a_i|)}
   1[a_j|\cdots|a_n|a_0|\cdots|a_{j-1}].
\end{align*}

In the $j=0$ term of $b$, the product is the outer
coefficient. A scalar in a bar slot gives zero. Adjacent
mergers cancel in pairs, as do the cyclic faces, giving
$d_{\mathrm H}^2=B^2=d_{\mathrm H}B+Bd_{\mathrm H}=0$.

Let $\zeta$ have degree two. The negative-cyclic complex is
\[
 M_R=\bigl(C_\bullet(R)[[\zeta]],
                  d_{\mathrm H}+\zeta B\bigr).
\]
Its volume cycle is
\begin{equation}\label{mht:eq:resolved-volume-cycle}
 \omega_0=1[q|p]-1[p|q]-1[t].
\end{equation}

The Hochschild boundary of the first two terms is
$-1[qp-pq]$. In this convention the internal differential
of $-1[t]$ is $+1[qp-pq]$. Every Connes term has a unit in
a bar slot. Hence $d_{\mathrm H}\omega_0=B\omega_0=0$.

Put $\Omega_A^a$ in degree $-a$. The normalized
Hochschild--Kostant--Rosenberg projection is
\begin{equation}\label{mht:eq:resolved-hkr-projection}
 P(a_0[a_1|\cdots|a_n])
 =\frac1{n!}\pi(a_0)\,
       d_{\mathrm{dR}}\pi(a_1)\wedge\cdots\wedge
       d_{\mathrm{dR}}\pi(a_n).
\end{equation}
It satisfies $Pd_{\mathrm H}=0$, $PB=d_{\mathrm{dR}}P$
and $P\omega_0=dq\wedge dp$.

Give $q,p,dq,dp$ weight one, $t$ weight two, and $\zeta$
weight zero. Every normalized bar slot has positive weight.
At fixed weight, bar length and internal degree are bounded,
so each chain space is finite dimensional.

\begin{proposition}\label{mht:prop:mixed-volume-contraction}
There is a weight-preserving contraction $(P,I,T)$ from
$(C_\bullet(R),d_{\mathrm H})$ to $(\Omega_A^\bullet,0)$,
with $I(dq\wedge dp)=\omega_0$. It extends to a contraction
from $M_R$ to $(\Omega_A^\bullet[[\zeta]],\zeta d_{\mathrm{dR}})$,
with unchanged projection and
\begin{equation}\label{mht:eq:mixed-perturbation}
 I_\zeta=(1+T\zeta B)^{-1}I,
 \qquad T_\zeta=T(1+\zeta BT)^{-1}.
\end{equation}
The inclusion still sends $dq\wedge dp$ to $\omega_0$.
\end{proposition}

\begin{proof}
The contraction of $R\to A$ induces a quasi-isomorphism of
Hochschild chains. Indeed the bar-length filtration is finite
in each weight, and its associated-graded comparison is the
tensor product of quasi-isomorphisms over $\mathbb C$.

For $A$, the bimodule Koszul complex on
$q_L-q_R,p_L-p_R$ resolves the diagonal. The first element
is a non-zero-divisor in $A^e$. The second remains one in
its quotient. Successive two-term resolutions prove exactness.

Tensoring with the diagonal kills both differentials and
gives $\Omega_A^\bullet$. Antisymmetrization identifies the
projection with \eqref{mht:eq:resolved-hkr-projection}.
The normalized formula for $B$ becomes $d_{\mathrm{dR}}$.

Thus $P$ is a surjective quasi-isomorphism at each weight.
Lift a basis of forms to cycles, choosing $\omega_0$ for the
volume. Split its acyclic kernel in each degree as
\[
 (\ker P)^r=B^r\oplus L^r,\qquad
 d_{\mathrm H}:L^r\xrightarrow{\sim}B^{r+1}.
\]
Let $T$ be this inverse on $B^{r+1}$ and zero on $L$ and
the chosen lifts. Then
$d_{\mathrm H}T+Td_{\mathrm H}=1-IP$ and
$T^2=TI=PT=0$. These are finite linear-algebra choices.

Perturb by $\zeta B$. Since $P\zeta BT
=\zeta d_{\mathrm{dR}}PT=0$, the projection stays $P$ and
the transferred differential has no higher terms beyond
$\zeta d_{\mathrm{dR}}$. The series in
\eqref{mht:eq:mixed-perturbation} converge in $\zeta$.
At fixed weight and total degree they are finite, because
bar length and internal degree bound the possible powers
of $\zeta$. Finally $B\omega_0=0$ keeps its inclusion fixed.
\end{proof}

\section{Coherent invariance of the volume}

For a degree-$r$ derivation $D$, its action on raw chains is
\begin{equation}\label{mht:eq:raw-lie-derivative}
 \mathcal L_D(a_0[a_1|\cdots|a_n])
 =\sum_{j=0}^n(-1)^{nr+r\sum_{i<j}|a_i|}
    a_0[\cdots|D(a_j)|\cdots|a_n].
\end{equation}
The $j=0$ term acts on the outer coefficient. The factor
$(-1)^{nr}$ matches the internal differential in
\eqref{mht:eq:raw-mixed-convention}. Tensor signs give
\[
 [d_{\mathrm H},\mathcal L_D]=\mathcal L_{[d,D]},
 \qquad [B,\mathcal L_D]=0,
 \qquad [\mathcal L_D,\mathcal L_E]=\mathcal L_{[D,E]}.
\]
The action $F$ therefore acts on the negative-cyclic complex.

More explicitly, regard $F_n$ as the values of a degree-one
twisting cochain
$\alpha:S^c(\mathfrak l[1])\to\operatorname{Der}(R)$.
Its equation, with the tensor sign in the convolution
bracket, is
\[
 \partial\alpha+\alpha b_{\mathfrak l}
       +\tfrac12[\alpha,\alpha]=0.
\]
Its arity-two equation is exactly
$[d,F_2(f,g)]=[F_1(f),F_1(g)]-F_1\{f,g\}$.

Use the arity product of maps from this coaugmented
coalgebra to $M_R$. For a homogeneous map $\phi$, its
differential is
\[
 (d_{\mathrm H}+\zeta B)\phi
 -(-1)^{|\phi|}\phi b_{\mathfrak l}
 +\mathcal L_\alpha\star\phi.
\]
The last product uses the unshuffle coproduct and its tensor
Koszul sign. Let $\delta_F$ denote the terms beyond the
internal differential. They raise arity.

The projection $P$ intertwines $F_1$ with $X_f$ on forms.
It kills the action of every $F_n$ for $n\geq2$, since
their negative-degree generator values are killed by $\pi$.
Thus it intertwines $\delta_F$ with the strict Hamiltonian
action on the corresponding convolution of forms.

\begin{theorem}\label{mht:thm:coherent-volume}
For the fixed classical action $F$, the volume $dq\wedge dp$
has a coherent negative-cyclic lift
\begin{equation}\label{mht:eq:coherent-volume}
 \omega_{\mathrm{eq}}
 =(1+\mathcal T_\zeta\delta_F)^{-1}
                  I_\zeta(dq\wedge dp),
\end{equation}
where $\mathcal T_\zeta$ is the contraction on convolution
induced by $T_\zeta$. Its arity-$n$ component has values
in $M_R^{-n-2}$. The derived fiber over this specified
invariant volume is contractible. The construction extends
continuously to pointed formal Hamiltonians at fixed arity.
\end{theorem}

\begin{proof}
Extend Proposition~\ref{mht:prop:mixed-volume-contraction}
to convolution and perturb by $\delta_F$. This series
converges in arity. Since the projection intertwines the
actions and $P\mathcal T_\zeta=0$, one has
$P\delta_F\mathcal T_\zeta=0$. The transferred differential
therefore consists of $\zeta d_{\mathrm{dR}}$ and the
strict Chevalley action on forms.

The volume is de Rham closed and Hamiltonian invariant:
$\operatorname{div}(X_f)=-f_{pq}+f_{qp}=0$.
Its perturbed inclusion is consequently closed. The same
perturbation contracts the projection kernel. The homotopy
fiber over the specified closed volume is an affine translate
of this acyclic complex, so its derived mapping space is
contractible. Both the action and its projection are fixed
in this assertion.

At fixed arity only finitely many arity perturbations
contribute. At fixed weight and degree the mixed perturbation
is finite. Proposition~\ref{mht:prop:pointed-action-completion}
gives nonnegative input weight shifts, leaving finitely many
input degrees at any fixed output weight. This proves
continuity for the pointed completion.
\end{proof}

The first two coherence equations have explicit representatives.
Write $\langle a_1|\cdots|a_n\rangle=1[a_1|\cdots|a_n]$.
The arity-one equation is
$d_{\mathrm H}\omega_1(f)=-\mathcal L_{F_1(f)}\omega_0$.
For cubics choose
\[
 \omega_1(q^3)=-3\langle q|q|q\rangle,
 \qquad \omega_1(p^3)=-3\langle p|p|p\rangle.
\]
A compatible quartic representative is
\begin{align}\label{mht:eq:quartic-volume-homotopy}
 \omega_1(q^2p^2)
 ={}&-\tfrac23\langle tt\rangle\\
 &+\tfrac23\bigl(
  \langle q|pt\rangle-\langle q|tp\rangle
 -\langle p|qt\rangle+\langle p|tq\rangle\bigr)\notag\\
 &+\tfrac23\bigl(
  \langle qt|p\rangle-\langle pt|q\rangle
 -\langle tq|p\rangle+\langle tp|q\rangle\bigr)\notag\\
 &-2\bigl(\langle q|p|qp\rangle
          +\langle p|qp|q\rangle
          +\langle qp|q|p\rangle\bigr).\notag
\end{align}
All these chains have zero $B$-image and zero HKR projection.

For $f=q^3,g=p^3$, the next equation in the same convention is
\begin{align}\label{mht:eq:second-volume-equation}
 d_{\mathrm H}\omega_2(f,g)
 ={}&\omega_1(\{f,g\})
 -\mathcal L_{F_1(f)}\omega_1(g)
 +\mathcal L_{F_1(g)}\omega_1(f)\notag\\
 &-\mathcal L_{F_2(f,g)}\omega_0.
\end{align}
The minus sign uses $F_2=\eta$ as in
\eqref{mht:eq:cubic-action-homotopy}. One solution is
\begin{align}\label{mht:eq:second-volume-homotopy}
 \tfrac19\omega_2(q^3,p^3)
 ={}&\langle t|t\rangle
 +\langle q|p|t\rangle-\langle p|q|t\rangle
 -\langle q|t|p\rangle+\langle p|t|q\rangle\notag\\
 &+\langle t|q|p\rangle-\langle t|p|q\rangle
 -\langle q|q|p|p\rangle+\langle q|p|q|p\rangle\notag\\
 &+\langle q|p|p|q\rangle+\langle p|q|q|p\rangle
 -\langle p|q|p|q\rangle-\langle p|p|q|q\rangle.
\end{align}
Its Connes image and HKR projection vanish. Substitution in
the raw differential verifies both displayed coherence
equations: the internal terms use $dt=qp-pq$, and the
adjacent mergers cancel the cyclic faces. The terms
$\langle tt\rangle$ and $\langle t|t\rangle$ belong to
different equations and are both necessary.

These prescribed lower components extend to all arities
because their later residuals lie in the contracted projection
kernel. They need not coincide with the representatives
selected by a different finite row-reduction convention.

The volume also records nondegeneracy. The diagonal Koszul
resolution of $A$ is finite free over $A^e$ and self-dual,
with its top generator paired with $dq\wedge dp$. Hence
\[
 R\operatorname{Hom}_{A^e}(A,A^e)\simeq A[-2].
\]
Transport through $R\to A$ gives the smooth
$2$-Calabi--Yau orientation of the associative resolution.
This duality uses the diagonal bimodule. A proper scalar
trace would require a different finiteness condition.

\section{The quantum cycle and the remaining lift}

For $\mathcal R_\varepsilon$, take normalized chains relative
to $k_\varepsilon$ with the same raw convention. The expression
$1[q|p]-1[p|q]-1[t]$ is still a negative-cyclic cycle.
Its differential is $-\varepsilon1[1]$, which vanishes in
the normalized complex, and its Connes image is zero.

This produces the quantum volume cycle and the coherent
quantum derivation action separately. Their compatible
equivariant cyclic comparison requires its own construction.
On the classical matrix gauge complex, a Hamiltonian-function
lift must retain the central rank or specify the pointed source.
Its moment-map equations still require compatible solutions.
The associative derivation contraction supplies neither
conclusion by forgetting the extra structure.
